%%%%%%%%%%%%%%%%%%%%%%%%%%%%%%%%%%%%%%%%%%%%%%%%%%%%%%%%%%%%%%%%%%%
% Physik 4 - Blatt 09 Abgabe
%%%%%%%%%%%%%%%%%%%%%%%%%%%%%%%%%%%%%%%%%%%%%%%%%%%%%%%%%%%%%%%%%%%

\documentclass[11pt,a4paper]{scrartcl}

% Layout
\usepackage[margin=2.5cm]{geometry}
\usepackage[onehalfspacing]{setspace}

% General packages
\usepackage{graphicx}
\usepackage{enumitem}
\usepackage{tcolorbox}
\tcbuselibrary{breakable}
\usepackage[breaklinks=true,colorlinks=true,linkcolor=blue,urlcolor=blue,citecolor=blue]{hyperref}
\usepackage{amsmath,amsthm,amssymb,mathtools}
\usepackage{icomma}
\usepackage[T1]{fontenc}
\usepackage[utf8]{inputenc}
\usepackage[ngerman]{babel}
\usepackage{tikz}

% Units / tables / floats / references
\usepackage[locale = DE,space-before-unit=true,per-mode = symbol]{siunitx}
\usepackage{booktabs,multirow}
\usepackage{placeins}
\usepackage[natbib,abbreviate=true,doi=false,style=numeric-comp,giveninits=true,sorting=none]{biblatex}
\usepackage{csquotes}
\usepackage{pgfplots}
\pgfplotsset{compat=1.18}

% Graphics paths
\graphicspath{{Bilder/} {images/} {figures/} {chapters/}}

% Optional math shortcuts
\newcommand{\R}{\mathbb{R}}
\newcommand{\N}{\mathbb{N}}
\newcommand{\Q}{\mathbb{Q}}
\newcommand{\set}[1]{\left\{ #1 \right\}}
\newcommand{\abs}[1]{\left| #1 \right|}
\newcommand{\norm}[1]{\left\lVert #1 \right\rVert}
\newcommand{\ol}[1]{\overline{#1}}
\newcommand{\half}{\frac{1}{2}}
\newcommand{\dd}{\,\mathrm{d}}
\newcommand{\pifour}{p_{\pi}}

% Optional units
\DeclareSIUnit{\year}{a}
\DeclareSIUnit{\day}{d}

% Box styles
\newtcolorbox{calculationbox}[1][]{
    title=Calculation,
    breakable,
    #1
}

\newtcolorbox{solutionbox}[1][]{
    title=Solution,
    breakable,
    #1
}

\newtcolorbox{answerbox}[1][]{
    title=Answer,
    breakable,
    #1
}

\newtcolorbox{notebox}[1][]{
    title=Notes,
    breakable,
    #1
}

\newtcolorbox{sidecalcbox}[1][]{
    title=Side Calculation,
    breakable,
    #1
}

%%%%%%%%%%%%%%%%%%%%%%%%%%%%%%%%%%%%%%%%%%%%%%%%%%%%%%%%%%%%%%%%%%%
\begin{document}

    \title{Physik 4 - Blatt 09}
    \author{Tobias Laser}
    \date{\today}
    \maketitle
    \vfill
    \thispagestyle{empty}
    \newpage

    \pagenumbering{arabic}

    \paragraph{Aufgabe 1: Alles ist relativ (egal)}

    Betrachten Sie die Reaktion
    \[
        \pi^- + p \longrightarrow \Delta^0 \longrightarrow \pi^0+n.
    \]
    Nehmen Sie dabei an, dass das Proton in Ruhe sein soll. Rechnen Sie im Folgenden
    in Einheiten von MeV mit einer Genauigkeit von zwei Nachkommastellen.

    \begin{enumerate}[label=(\alph*)]
        \item
            Berechnen Sie als Erstes die Schwellenergie des $\pi^-$ für die Erzeugung der
            $\Delta^0$-Resonanz, indem Sie allgemein die Schwerpunktsenergie bestimmen
            und diese als die Masse der $\Delta^0$-Resonanz annehmen.

            \begin{calculationbox}
                Die Idee ist, mit einer Lorentz-invarianten Größe zu arbeiten. Dafür verwendet
                man das Quadrat der Gesamtviererimpulses:
                \[
                    s = (p_{\pi}+p_p)^2.
                \]
                Im Schwerpunktsystem ist die Gesamtimpuls-Summe null. Daher gilt dort
                \[
                    \sqrt{s}=E_{\mathrm{cm}}.
                \]
                An der Schwelle soll gerade die $\Delta^0$-Resonanz erzeugt werden. Also muss
                im Schwerpunktsystem gerade ihre Ruheenergie verfügbar sein:
                \[
                    E_{\mathrm{cm}}=m_{\Delta}c^2.
                \]

                \begin{calculationbox}
                    Wir rechnen in Teilchenphysik-Einheiten mit $c=1$. Dann haben Massen die
                    Einheit MeV und Impulse die Einheit MeV. Im Laborsystem ruht das Proton:
                    \[
                        p_p=(m_p,0),
                        \qquad
                        p_{\pi}=(E_{\pi},\vec p_{\pi}).
                    \]
                    Damit ist
                    \begin{align*}
                        s
                        &= (p_{\pi}+p_p)^2\\
                        &= p_{\pi}^2+p_p^2+2p_{\pi}\cdot p_p\\
                        &= m_{\pi}^2+m_p^2+2m_pE_{\pi}.
                    \end{align*}
                    An der Schwelle setzen wir $s=m_{\Delta}^2$:
                    \[
                        m_{\Delta}^2=m_{\pi}^2+m_p^2+2m_pE_{\pi,\mathrm{thr}}.
                    \]
                    Auflösen nach der Pionenergie liefert
                    \[
                        E_{\pi,\mathrm{thr}}
                        =\frac{m_{\Delta}^2-m_p^2-m_{\pi}^2}{2m_p}.
                    \]
                \end{calculationbox}

                \begin{calculationbox}
                    Mit
                    \[
                        m_{\Delta}=\SI{1232}{MeV},
                        \qquad
                        m_p=\SI{938.3}{MeV},
                        \qquad
                        m_{\pi^-}=\SI{139.6}{MeV}
                    \]
                    folgt
                    \begin{align*}
                        E_{\pi,\mathrm{thr}}
                        &=\frac{(1232)^2-(938.3)^2-(139.6)^2}{2\cdot 938.3}\,\si{MeV}\\
                        &=\SI{329.28}{MeV}.
                    \end{align*}
                    Falls man die kinetische Energie des Pions meint, muss noch die Ruheenergie
                    des Pions abgezogen werden:
                    \begin{align*}
                        T_{\pi,\mathrm{thr}}
                        &=E_{\pi,\mathrm{thr}}-m_{\pi}\\
                        &=\SI{329.28}{MeV}-\SI{139.60}{MeV}\\
                        &=\SI{189.68}{MeV}.
                    \end{align*}
                \end{calculationbox}
            \end{calculationbox}

            \begin{answerbox}
                Die Schwellenergie als Gesamtenergie des einlaufenden Pions ist
                \[
                    \boxed{E_{\pi,\mathrm{thr}}=\SI{329.28}{MeV}.}
                \]
                Die dazugehörige kinetische Schwellenergie ist
                \[
                    \boxed{T_{\pi,\mathrm{thr}}=\SI{189.68}{MeV}.}
                \]
                Für die Rechnung im nächsten Teil verwenden wir die Gesamtenergie
                $E_{\pi,\mathrm{thr}}$.
            \end{answerbox}

        \item
            Bestimmen Sie für die Situation, dass das Pion gerade die Schwellenergie hat,
            die Geschwindigkeit in Einheiten der Lichtgeschwindigkeit, mit der sich das
            Schwerpunktsystem relativ zum Laborsystem bewegt.

            \begin{calculationbox}
                Das Schwerpunktsystem bewegt sich relativ zum Laborsystem mit der Geschwindigkeit,
                bei der der Gesamtimpuls verschwindet. Für ein System aus Pion und ruhendem Proton gilt
                in Einheiten $c=1$:
                \[
                    \beta_{\mathrm{cm}}
                    =\frac{\abs{\vec p_{\mathrm{ges}}}}{E_{\mathrm{ges}}}
                    =\frac{p_{\pi}}{E_{\pi}+m_p}.
                \]
                Der Impuls des Pions bei Schwellenergie folgt aus der relativistischen
                Energie-Impuls-Beziehung:
                \[
                    E_{\pi}^2=p_{\pi}^2+m_{\pi}^2.
                \]

                \begin{calculationbox}
                    Also
                    \begin{align*}
                        p_{\pi,\mathrm{thr}}
                        &=\sqrt{E_{\pi,\mathrm{thr}}^2-m_{\pi}^2}\\
                        &=\sqrt{(329.28)^2-(139.6)^2}\,\si{MeV}\\
                        &=\SI{298.22}{MeV}.
                    \end{align*}
                    Damit ist
                    \begin{align*}
                        \beta_{\mathrm{cm}}
                        &=\frac{298.22}{329.28+938.3}\\
                        &=0.23527.
                    \end{align*}
                \end{calculationbox}
            \end{calculationbox}

            \begin{answerbox}
                Das Schwerpunktsystem bewegt sich relativ zum Laborsystem mit
                \[
                    \boxed{\beta_{\mathrm{cm}}=0.24}
                \]
                beziehungsweise genauer mit
                \[
                    \boxed{v_{\mathrm{cm}}=0.23527\,c.}
                \]
            \end{answerbox}

        \item
            Im Schwerpunktsystem zerfällt das $\Delta^0$ nach kurzer Zeit in ein neutrales Pion
            und ein Neutron. Dafür kann man berechnen, dass die Energie des Neutrons im
            Schwerpunktsystem \SI{966.88}{MeV} und sein Impuls \SI{228.19}{MeV/c} beträgt.
            Bestimmen Sie daraus die maximale und die minimale Energie, die das Neutron
            im Laborsystem haben kann.

            \begin{calculationbox}
                Wir transformieren die Neutronenergie aus dem Schwerpunktsystem in das Laborsystem.
                Die Lorentztransformation für die Energie lautet bei einem Boost entlang der
                Bewegungsrichtung des Schwerpunktsystems
                \[
                    E_{\mathrm{lab}}=\gamma_{\mathrm{cm}}
                    \left(E^*+\beta_{\mathrm{cm}}p^*\cos\theta^*\right).
                \]
                Dabei bezeichnet der Stern Größen im Schwerpunktsystem. Der Winkel $\theta^*$ ist
                der Winkel zwischen der Neutron-Flugrichtung im Schwerpunktsystem und der Boostrichtung.

                \begin{calculationbox}
                    Zuerst berechnen wir
                    \[
                        \gamma_{\mathrm{cm} }
                        =\frac{1}{\sqrt{1-\beta_{\mathrm{cm}}^2}}
                        =\frac{1}{\sqrt{1-0.23527^2}}
                        =1.02888.
                    \]
                    Die maximale Energie entsteht, wenn das Neutron im Schwerpunktsystem in
                    Boostrichtung fliegt, also $\cos\theta^*=+1$:
                    \begin{align*}
                        E_{\max}
                        &=\gamma_{\mathrm{cm}}(E_n^*+\beta_{\mathrm{cm}}p_n^*)\\
                        &=1.02888\left(966.88+0.23527\cdot 228.19\right)\,\si{MeV}\\
                        &=\SI{1050.04}{MeV}.
                    \end{align*}
                    Die minimale Energie entsteht, wenn das Neutron entgegen der Boostrichtung
                    fliegt, also $\cos\theta^*=-1$:
                    \begin{align*}
                        E_{\min}
                        &=\gamma_{\mathrm{cm}}(E_n^*-\beta_{\mathrm{cm}}p_n^*)\\
                        &=1.02888\left(966.88-0.23527\cdot 228.19\right)\,\si{MeV}\\
                        &=\SI{939.57}{MeV}.
                    \end{align*}
                \end{calculationbox}
            \end{calculationbox}

            \begin{answerbox}
                Die maximale Neutronenergie im Laborsystem ist
                \[
                    \boxed{E_{\max}=\SI{1050.04}{MeV}.}
                \]
                Die minimale Neutronenergie im Laborsystem ist
                \[
                    \boxed{E_{\min}=\SI{939.57}{MeV}.}
                \]
                Interpretation: Vorwärts emittierte Neutronen werden durch den Boost energetisch
                angehoben, rückwärts emittierte Neutronen energetisch abgesenkt.
            \end{answerbox}
    \end{enumerate}

    \newpage
    \paragraph{Aufgabe 2: Delta und (invariante) Masse}

    Ein $\Delta^0$-Baryon mit
    \[
        m_{\Delta^0}=\SI{1232}{MeV/c^2}
    \]
    und einer relativistischen Gesamtenergie von \SI{1.35}{GeV} fliegt im Laborsystem
    entlang der positiven $x$-Achse und zerfällt über die Reaktion
    \[
        \Delta^0 \longrightarrow p+\pi^-.
    \]
    Zusätzlich betrachten wir den Prozess in dem Schwerpunktsystem, welches sich durch
    eine einfache Lorentz-Transformation des Laborsystems ergibt. In diesem sollen die
    beiden Zerfallsprodukte entlang der $y$-Achse voneinander wegfliegen.

    \begin{enumerate}[label=(\alph*)]
        \item
            Machen Sie eine Skizze für beide Systeme.

            \begin{calculationbox}
                Im Laborsystem bewegt sich das $\Delta^0$ entlang der positiven $x$-Achse. Sein
                Zerfall findet während dieser Vorwärtsbewegung statt. Im Schwerpunktsystem ruht
                das $\Delta^0$ vor dem Zerfall. Dort fliegen Proton und Pion rein entlang der
                $y$-Achse auseinander. Wir wählen für die Rechnung
                \[
                    p \text{ in } +y\text{-Richtung},
                    \qquad
                    \pi^- \text{ in } -y\text{-Richtung}.
                \]

                \begin{center}
                    \begin{tikzpicture}[scale=1.0,>=stealth]
                        % left: lab
                        \draw[->] (-0.2,0) -- (4.2,0) node[right] {$x$};
                        \draw[->] (0,-1.2) -- (0,1.4) node[above] {$y$};
                        \node at (2,1.75) {Laborsystem};
                        \draw[->,thick] (0.4,0.15) -- (2.0,0.15) node[midway,above] {$\Delta^0$};
                        \fill (2.1,0.15) circle (1.5pt);
                        \draw[->,thick] (2.1,0.15) -- (3.4,0.85) node[right] {$p$};
                        \draw[->,thick] (2.1,0.15) -- (3.0,-0.75) node[right] {$\pi^-$};

                        % right: cm
                        \begin{scope}[xshift=7cm]
                            \draw[->] (-1.7,0) -- (1.7,0) node[right] {$x'$};
                            \draw[->] (0,-1.4) -- (0,1.4) node[above] {$y'$};
                            \node at (0,1.75) {Schwerpunktsystem};
                            \fill (0,0) circle (1.5pt) node[below right] {$\Delta^0$ ruht};
                            \draw[->,thick] (0,0) -- (0,1.0) node[right] {$p$};
                            \draw[->,thick] (0,0) -- (0,-1.0) node[right] {$\pi^-$};
                        \end{scope}
                    \end{tikzpicture}
                \end{center}
            \end{calculationbox}

            \begin{answerbox}
                Die Skizze zeigt den entscheidenden Unterschied: Im Laborsystem bewegt sich das
                Gesamtsystem nach rechts, im Schwerpunktsystem ruht das $\Delta^0$ und die beiden
                Zerfallsprodukte haben entgegengesetzte Impulse entlang der $y$-Achse.
            \end{answerbox}

        \item
            Mit welchem Bruchteil der Lichtgeschwindigkeit $\beta$ bewegt sich das $\Delta^0$
            im Laborsystem?

            \begin{calculationbox}
                Für ein einzelnes Teilchen gilt
                \[
                    E^2=p^2+m^2.
                \]
                Außerdem gilt für die Geschwindigkeit in Einheiten von $c$:
                \[
                    \beta=\frac{p}{E}.
                \]

                \begin{calculationbox}
                    Mit
                    \[
                        E_{\Delta}=\SI{1.35}{GeV},
                        \qquad
                        m_{\Delta}=\SI{1.232}{GeV}
                    \]
                    folgt zuerst
                    \begin{align*}
                        p_{\Delta}
                        &=\sqrt{E_{\Delta}^2-m_{\Delta}^2}\\
                        &=\sqrt{(1.35)^2-(1.232)^2}\,\si{GeV}\\
                        &=\SI{0.5520}{GeV}.
                    \end{align*}
                    Damit ist
                    \begin{align*}
                        \beta
                        &=\frac{p_{\Delta}}{E_{\Delta}}\\
                        &=\frac{0.5520}{1.35}\\
                        &=0.40887.
                    \end{align*}
                    Der Lorentzfaktor ist zugleich
                    \[
                        \gamma=\frac{E_{\Delta}}{m_{\Delta}}=\frac{1.35}{1.232}=1.09578.
                    \]
                \end{calculationbox}
            \end{calculationbox}

            \begin{answerbox}
                Das $\Delta^0$ bewegt sich im Laborsystem mit
                \[
                    \boxed{\beta=0.41}
                \]
                beziehungsweise genauer
                \[
                    \boxed{v=0.40887\,c.}
                \]
            \end{answerbox}

        \item
            Geben Sie die Viererimpulse von Proton und Pion im Schwerpunktsystem an in
            Einheiten von GeV/$c$ für alle Komponenten. Verwenden Sie dazu die Formel für
            die Aufteilung der Energie beim Zweikörperzerfall. Hinweis:
            $m_p=\SI{938.3}{MeV/c^2}$, $m_{\pi}=\SI{139.6}{MeV/c^2}$.

            \begin{calculationbox}
                Im Schwerpunktsystem ist das $\Delta^0$ in Ruhe. Daher ist die Gesamtenergie
                \[
                    E^*_{\mathrm{ges}}=m_{\Delta}.
                \]
                Beim Zweikörperzerfall $\Delta^0\to p+\pi^-$ sind die Energien im
                Schwerpunktsystem fest:
                \[
                    E_p^*=\frac{m_{\Delta}^2+m_p^2-m_{\pi}^2}{2m_{\Delta}},
                    \qquad
                    E_{\pi}^*=\frac{m_{\Delta}^2+m_{\pi}^2-m_p^2}{2m_{\Delta}}.
                \]

                \begin{calculationbox}
                    Wir verwenden GeV:
                    \[
                        m_{\Delta}=1.232,
                        \qquad
                        m_p=0.9383,
                        \qquad
                        m_{\pi}=0.1396.
                    \]
                    Damit
                    \begin{align*}
                        E_p^*
                        &=\frac{1.232^2+0.9383^2-0.1396^2}{2\cdot1.232}\,\si{GeV}\\
                        &=\SI{0.96540}{GeV},\\[0.5em]
                        E_{\pi}^*
                        &=\frac{1.232^2+0.1396^2-0.9383^2}{2\cdot1.232}\,\si{GeV}\\
                        &=\SI{0.26660}{GeV}.
                    \end{align*}
                    Der Impulsbetrag ist für beide Zerfallsprodukte gleich groß:
                    \begin{align*}
                        p^*
                        &=\sqrt{(E_p^*)^2-m_p^2}\\
                        &=\SI{0.22713}{GeV}.
                    \end{align*}
                \end{calculationbox}

                \begin{calculationbox}
                    Da die Teilchen entlang der $y$-Achse auseinanderfliegen, wählen wir
                    \[
                        \vec p_p^*=(0,+p^*,0),
                        \qquad
                        \vec p_{\pi}^*=(0,-p^*,0).
                    \]
                    In der Konvention
                    \[
                        p^{\mu}=\left(\frac{E}{c},p_x,p_y,p_z\right)
                    \]
                    erhalten wir in Einheiten von GeV/$c$:
                    \[
                        p_p^{*\mu}=(0.96540,0,+0.22713,0),
                    \]
                    \[
                        p_{\pi}^{*\mu}=(0.26660,0,-0.22713,0).
                    \]
                \end{calculationbox}
            \end{calculationbox}

            \begin{answerbox}
                Die Viererimpulse im Schwerpunktsystem sind
                \[
                    \boxed{p_p^{*\mu}=(0.96540,0,+0.22713,0)\,\si{GeV/c}}
                \]
                und
                \[
                    \boxed{p_{\pi}^{*\mu}=(0.26660,0,-0.22713,0)\,\si{GeV/c}.}
                \]
                Die Summe ist
                \[
                    p_p^{*\mu}+p_{\pi}^{*\mu}=(1.232,0,0,0)\,\si{GeV/c},
                \]
                also genau der Viererimpuls des ruhenden $\Delta^0$.
            \end{answerbox}

        \item
            Berechnen Sie dann die Viererimpulse von Proton und Pion im Laborsystem und
            bilden Sie zur Kontrolle die Summe. Hinweis: Benutzen Sie die Lorentztransformation
            und überlegen Sie sich das richtige Vorzeichen vor dem Diagonalelement $\gamma\beta$
            der Lorentzmatrix.

            \begin{calculationbox}
                Das Laborsystem erhält man aus dem Schwerpunktsystem durch einen Boost in positive
                $x$-Richtung. Für die Transformation vom Schwerpunktsystem ins Laborsystem gilt
                \[
                    \begin{pmatrix}
                        E\\ p_x\\ p_y\\ p_z
                    \end{pmatrix}_{\mathrm{lab}}
                    =
                    \begin{pmatrix}
                        \gamma & \gamma\beta & 0 & 0\\
                        \gamma\beta & \gamma & 0 & 0\\
                        0 & 0 & 1 & 0\\
                        0 & 0 & 0 & 1
                    \end{pmatrix}
                    \begin{pmatrix}
                        E^*\\ p_x^*\\ p_y^*\\ p_z^*
                    \end{pmatrix}.
                \]
                Hier ist das Pluszeichen richtig, weil wir vom Ruhesystem des $\Delta^0$ in das
                Laborsystem boosten, in dem das $\Delta^0$ in positive $x$-Richtung fliegt.

                \begin{calculationbox}
                    Da im Schwerpunktsystem $p_x^*=0$ gilt, vereinfacht sich die Transformation zu
                    \[
                        E_{\mathrm{lab}}=\gamma E^*,
                        \qquad
                        p_{x,\mathrm{lab}}=\gamma\beta E^*,
                        \qquad
                        p_{y,\mathrm{lab}}=p_y^*.
                    \]
                    Mit
                    \[
                        \beta=0.40887,
                        \qquad
                        \gamma=1.09578
                    \]
                    folgt für das Proton
                    \begin{align*}
                        E_p&=1.09578\cdot0.96540=\SI{1.05786}{GeV},\\
                        p_{p,x}&=1.09578\cdot0.40887\cdot0.96540=\SI{0.43253}{GeV},\\
                        p_{p,y}&=+\SI{0.22713}{GeV}.
                    \end{align*}
                    Also
                    \[
                        p_p^{\mu}=(1.05786,0.43253,+0.22713,0)\,\si{GeV/c}.
                    \]
                \end{calculationbox}

                \begin{calculationbox}
                    Für das Pion folgt analog
                    \begin{align*}
                        E_{\pi}&=1.09578\cdot0.26660=\SI{0.29214}{GeV},\\
                        p_{\pi,x}&=1.09578\cdot0.40887\cdot0.26660=\SI{0.11945}{GeV},\\
                        p_{\pi,y}&=-\SI{0.22713}{GeV}.
                    \end{align*}
                    Also
                    \[
                        p_{\pi}^{\mu}=(0.29214,0.11945,-0.22713,0)\,\si{GeV/c}.
                    \]
                \end{calculationbox}

                \begin{calculationbox}
                    Als Kontrolle bilden wir die Summe:
                    \begin{align*}
                        p_p^{\mu}+p_{\pi}^{\mu}
                        &=(1.05786+0.29214,\,0.43253+0.11945,\,0.22713-0.22713,\,0)\\
                        &=(1.35000,0.55198,0,0)\,\si{GeV/c}.
                    \end{align*}
                    Das passt zum Viererimpuls des $\Delta^0$ im Laborsystem:
                    \[
                        p_{\Delta}^{\mu}=(E_{\Delta},p_{\Delta},0,0)
                        =(1.35,0.55198,0,0)\,\si{GeV/c}.
                    \]
                \end{calculationbox}
            \end{calculationbox}

            \begin{answerbox}
                Die Viererimpulse im Laborsystem sind
                \[
                    \boxed{p_p^{\mu}=(1.05786,0.43253,+0.22713,0)\,\si{GeV/c}}
                \]
                und
                \[
                    \boxed{p_{\pi}^{\mu}=(0.29214,0.11945,-0.22713,0)\,\si{GeV/c}.}
                \]
                Die Summe ist
                \[
                    \boxed{p_p^{\mu}+p_{\pi}^{\mu}=(1.35000,0.55198,0,0)\,\si{GeV/c}.}
                \]
            \end{answerbox}

        \item
            Wie groß ist der Winkel $\alpha$ zwischen dem Impuls des Protons und der $x$-Achse
            im Laborsystem? Berechnen Sie analog den entsprechenden Winkel für das Pion.
            Wie groß ist also der Öffnungswinkel zwischen den Bahnen der beiden Zerfallsprodukte
            im Laborsystem?

            \begin{calculationbox}
                Der Winkel zur positiven $x$-Achse ergibt sich aus den Impulskomponenten:
                \[
                    \tan\alpha=\frac{\abs{p_y}}{p_x}.
                \]
                Für das Proton ist
                \begin{align*}
                    \alpha_p
                    &=\arctan\left(\frac{0.22713}{0.43253}\right)\\
                    &=\SI{27.70}{\degree}.
                \end{align*}
                Für das Pion ist
                \begin{align*}
                    \alpha_{\pi}
                    &=\arctan\left(\frac{0.22713}{0.11945}\right)\\
                    &=\SI{62.26}{\degree}.
                \end{align*}
                Da die beiden Teilchen auf verschiedenen Seiten der $x$-Achse fliegen, ist der
                Öffnungswinkel die Summe dieser beiden Winkel:
                \begin{align*}
                    \theta_{\mathrm{open}}
                    &=\alpha_p+\alpha_{\pi}\\
                    &=\SI{27.70}{\degree}+\SI{62.26}{\degree}\\
                    &=\SI{89.97}{\degree}.
                \end{align*}
            \end{calculationbox}

            \begin{answerbox}
                Der Winkel des Protons zur $x$-Achse ist
                \[
                    \boxed{\alpha_p=\SI{27.70}{\degree}.}
                \]
                Der entsprechende Winkel des Pions zur $x$-Achse ist
                \[
                    \boxed{\alpha_{\pi}=\SI{62.26}{\degree}.}
                \]
                Der Öffnungswinkel zwischen beiden Bahnen ist
                \[
                    \boxed{\theta_{\mathrm{open}}=\SI{89.97}{\degree}.}
                \]
            \end{answerbox}

        \item
            Welche Geschwindigkeit ist größer: Die des Protons oder die des Pions? Versuchen
            Sie ohne Rechnung eine Antwort zu finden, und zwar zuerst für das Schwerpunktsystem
            und dann für das Laborsystem.

            \begin{calculationbox}
                Im Schwerpunktsystem haben Proton und Pion betragsmäßig denselben Impuls, weil
                sie aus dem ruhenden $\Delta^0$ in entgegengesetzte Richtungen auseinanderfliegen.
                Die Geschwindigkeit eines relativistischen Teilchens ist
                \[
                    v=\frac{pc^2}{E}.
                \]
                Bei gleichem Impuls hat das leichtere Teilchen die kleinere Energie und deshalb
                die größere Geschwindigkeit. Da das Pion viel leichter ist als das Proton, ist im
                Schwerpunktsystem das Pion schneller.

                \begin{calculationbox}
                    Auch im Laborsystem bleibt das Pion schneller. Man kann das qualitativ so sehen:
                    Der Boost in $x$-Richtung gibt beiden Teilchen dieselbe zusätzliche
                    $x$-Geschwindigkeit des Schwerpunktsystems,
                    \[
                        v_x=\beta c.
                    \]
                    Die transversale Geschwindigkeitskomponente wird aber durch den Faktor
                    $\gamma$ und durch die jeweilige Energie skaliert. Das Pion hatte im
                    Schwerpunktsystem die größere Geschwindigkeit und behält auch im Laborsystem
                    den größeren Geschwindigkeitsbetrag.
                \end{calculationbox}

                \begin{calculationbox}
                    Als Kontrolle kann man die Beträge berechnen:
                    \begin{align*}
                        \frac{v_p}{c}
                        &=\frac{\sqrt{0.43253^2+0.22713^2}}{1.05786}
                        =0.46182,\\
                        \frac{v_{\pi}}{c}
                        &=\frac{\sqrt{0.11945^2+0.22713^2}}{0.29214}
                        =0.87844.
                    \end{align*}
                \end{calculationbox}
            \end{calculationbox}

            \begin{answerbox}
                Sowohl im Schwerpunktsystem als auch im Laborsystem ist das Pion schneller:
                \[
                    \boxed{v_{\pi}>v_p.}
                \]
                Im Laborsystem erhält man zur Kontrolle ungefähr
                \[
                    \boxed{v_p=0.46c,\qquad v_{\pi}=0.88c.}
                \]
            \end{answerbox}
    \end{enumerate}

\end{document}
